\chapter{Pipelines et automatisation}
\label{ch:pipelines}

\begin{quote}
\emph{<< Si vous le faites deux fois, automatisez-le. >>}
--- Tous les ing\'enieurs logiciels
\end{quote}

% --- hook ---
La fuite de données la plus pernicieuse n'est pas un bug : c'est une convention qui paraît raisonnable. Vous calez votre scaler sur le jeu complet pour <<~économiser une étape~>> --- l'information du test infiltre alors la moyenne et l'écart-type du scaler. Vous validez par CV un modèle qui a déjà << vu~>> les valeurs qu'on prétend lui cacher. Votre score d'évaluation est artificiellement élevé, et vous ne le saurez qu'en production, quand le modèle réel se comportera moins bien que celui de votre rapport. Les \texttt{Pipeline} et \texttt{ColumnTransformer} de scikit-learn empaquettent le workflow une fois pour toutes et garantissent qu'aucune information ne traverse la cloison entre train et test. Ce chapitre apprend à les utiliser comme la discipline qu'ils incarnent, pas comme une API.
\section{Pourquoi des pipelines ?}

Le pr\'etraitement manuel est sujet aux erreurs et aux fuites. Erreurs courantes :
\begin{itemize}
  \item Ajuster un scaler sur tout le dataset (y compris le test) --- fuite de donn\'ees.
  \item Oublier d'appliquer les m\^emes transformations \`a la pr\'ediction.
  \item Appliquer les transformations dans un ordre diff\'erent entre train et inference.
\end{itemize}

Les \texttt{Pipeline} et \texttt{ColumnTransformer} de scikit-learn r\'esolvent les trois en empaquetant tout le workflow dans un objet unique et reproductible.

\begin{datatip}
Un pipeline garantit que chaque transformation appliqu\'ee \`a l'entra\^inement est appliqu\'ee identiquement \`a la pr\'ediction, dans le m\^eme ordre, avec les m\^emes param\`etres ajust\'es.
\end{datatip}

% ============================================================
\section{Pipeline basique}

\begin{minted}{python}
import pandas as pd
import numpy as np
from sklearn.pipeline import Pipeline
from sklearn.preprocessing import StandardScaler
from sklearn.impute import SimpleImputer
from sklearn.linear_model import LogisticRegression
from sklearn.model_selection import train_test_split

# Adult Census
url = ("https://archive.ics.uci.edu/ml/machine-learning-databases/"
       "adult/adult.data")
cols = ["age", "workclass", "fnlwgt", "education", "education_num",
        "marital_status", "occupation", "relationship", "race",
        "sex", "capital_gain", "capital_loss", "hours_per_week",
        "native_country", "income"]
df = pd.read_csv(url, header=None, names=cols, na_values=" ?",
                 skipinitialspace=True)

# Pipeline num\'erique simple
num_pipe = Pipeline([
    ("imputer", SimpleImputer(strategy="median")),
    ("scaler", StandardScaler()),
])

# Fit et transform
num_cols = ["age", "education_num", "capital_gain",
            "capital_loss", "hours_per_week"]
X_num = num_pipe.fit_transform(df[num_cols])
print(f"Shape: {X_num.shape}")
print(f"Means (should be ~0): {X_num.mean(axis=0).round(2)}")
\end{minted}

% ============================================================
\section{ColumnTransformer}

\begin{minted}{python}
from sklearn.compose import ColumnTransformer
from sklearn.preprocessing import OneHotEncoder

# D\'efinir les groupes de colonnes
num_features = ["age", "education_num", "capital_gain",
                "capital_loss", "hours_per_week"]
cat_features = ["workclass", "marital_status", "occupation",
                "relationship", "race", "sex"]

# Construire les transformateurs par type
num_transformer = Pipeline([
    ("imputer", SimpleImputer(strategy="median")),
    ("scaler", StandardScaler()),
])

cat_transformer = Pipeline([
    ("imputer", SimpleImputer(strategy="most_frequent")),
    ("encoder", OneHotEncoder(handle_unknown="ignore",
                               sparse_output=False)),
])

# Combiner avec ColumnTransformer
preprocessor = ColumnTransformer([
    ("num", num_transformer, num_features),
    ("cat", cat_transformer, cat_features),
])

# Fit et transform
X = preprocessor.fit_transform(df)
print(f"Input shape: {df[num_features + cat_features].shape}")
print(f"Output shape: {X.shape}")
\end{minted}

\begin{preprocesstip}
Utilisez \texttt{ColumnTransformer} pour appliquer diff\'erentes transformations \`a diff\'erents types de colonnes en une seule \'etape. C'est l'approche standard en production.
\end{preprocesstip}

% ============================================================
\section{Pipeline complet avec un mod\`ele}

\begin{minted}{python}
from sklearn.metrics import accuracy_score, classification_report

# Pr\'eparer la cible
y = (df["income"] == ">50K").astype(int)

# Pipeline complet : pr\'etraitement + mod\`ele
full_pipeline = Pipeline([
    ("preprocessor", preprocessor),
    ("classifier", LogisticRegression(max_iter=1000, random_state=42)),
])

# Split train/test
X_train, X_test, y_train, y_test = train_test_split(
    df[num_features + cat_features], y,
    test_size=0.2, random_state=42, stratify=y
)

# Ajuster tout le pipeline
full_pipeline.fit(X_train, y_train)

# Pr\'edire
y_pred = full_pipeline.predict(X_test)
print(f"Accuracy: {accuracy_score(y_test, y_pred):.4f}")
print(classification_report(y_test, y_pred))
\end{minted}

\begin{minted}{python}
# Validation crois\'ee avec le pipeline (pas de fuite !)
from sklearn.model_selection import cross_val_score

scores = cross_val_score(full_pipeline,
                         df[num_features + cat_features], y,
                         cv=5, scoring="accuracy")
print(f"CV Accuracy: {scores.mean():.4f} +/- {scores.std():.4f}")
\end{minted}

\begin{warningbox}
Si vous mettez \`a l'\'echelle vos donn\'ees \emph{avant} le split train/test, l'information du test fuite vers la moyenne et l'\'ecart-type du scaler. En pla\c{c}ant le scaler dans un pipeline, scikit-learn garantit qu'il ne s'ajuste que sur l'entra\^inement pendant la CV.
\end{warningbox}

% ============================================================
\section{Transformateurs personnalis\'es}

\begin{minted}{python}
from sklearn.base import BaseEstimator, TransformerMixin

class MissingIndicator(BaseEstimator, TransformerMixin):
    """Add binary columns indicating which values were missing."""

    def fit(self, X, y=None):
        self.columns_ = X.columns if hasattr(X, "columns") else None
        return self

    def transform(self, X):
        X = pd.DataFrame(X, columns=self.columns_)
        indicators = X.isnull().astype(int)
        indicators.columns = [f"{c}_missing" for c in indicators.columns]
        return pd.concat([X, indicators], axis=1)


class LogTransformer(BaseEstimator, TransformerMixin):
    """Apply log1p transformation to specified columns."""

    def __init__(self, columns=None):
        self.columns = columns

    def fit(self, X, y=None):
        return self

    def transform(self, X):
        X = X.copy()
        cols = self.columns or X.columns
        for col in cols:
            X[col] = np.log1p(X[col].clip(lower=0))
        return X
\end{minted}

\begin{minted}{python}
# Utiliser des transformateurs personnalis\'es dans un pipeline
custom_pipe = Pipeline([
    ("missing_ind", MissingIndicator()),
    ("imputer", SimpleImputer(strategy="median")),
    ("log", LogTransformer(columns=["capital_gain", "capital_loss"])),
    ("scaler", StandardScaler()),
])

# Tester
sample = df[num_features].head(20)
result = custom_pipe.fit_transform(sample)
print(f"Input columns: {len(num_features)}")
print(f"Output columns: {result.shape[1]}")
\end{minted}

% ============================================================
\section{Sauvegarder et charger un pipeline}

\begin{minted}{python}
import joblib

# Sauvegarder le pipeline ajust\'e
joblib.dump(full_pipeline, "income_pipeline.joblib")
print("Pipeline saved.")

# Le recharger plus tard (web app ou batch)
loaded_pipeline = joblib.load("income_pipeline.joblib")

# Pr\'edire sur de nouvelles donn\'ees
new_data = pd.DataFrame({
    "age": [35], "education_num": [13],
    "capital_gain": [0], "capital_loss": [0],
    "hours_per_week": [40], "workclass": ["Private"],
    "marital_status": ["Married-civ-spouse"],
    "occupation": ["Exec-managerial"],
    "relationship": ["Husband"], "race": ["White"],
    "sex": ["Male"],
})

prediction = loaded_pipeline.predict(new_data)
print(f"Prediction: {'> 50K' if prediction[0] else '<= 50K'}")
\end{minted}

\begin{preprocesstip}
Un pipeline sauvegard\'e contient tous les param\`etres ajust\'es (moyennes du scaler, cat\'egories de l'encodeur, statistiques d'imputation). Ce seul fichier suffit pour pr\'etraiter les nouvelles donn\'ees identiquement aux donn\'ees d'entra\^inement.
\end{preprocesstip}

% ============================================================
\section{Pipeline complet Melbourne Housing}

\begin{minted}{python}
# Pipeline complet pour Melbourne Housing
melb = pd.read_csv("melb_data.csv")
melb = melb.dropna(subset=["Price"])

num_features_melb = ["Rooms", "Distance", "Landsize", "BuildingArea",
                     "YearBuilt", "Car", "Bathroom"]
cat_features_melb = ["Type", "Method", "Regionname"]

melb_preprocessor = ColumnTransformer([
    ("num", Pipeline([
        ("imputer", SimpleImputer(strategy="median")),
        ("scaler", StandardScaler()),
    ]), num_features_melb),
    ("cat", Pipeline([
        ("imputer", SimpleImputer(strategy="most_frequent")),
        ("encoder", OneHotEncoder(handle_unknown="ignore",
                                   sparse_output=False)),
    ]), cat_features_melb),
])

from sklearn.ensemble import RandomForestRegressor

melb_pipeline = Pipeline([
    ("preprocessor", melb_preprocessor),
    ("model", RandomForestRegressor(n_estimators=100, random_state=42)),
])

X_melb = melb[num_features_melb + cat_features_melb]
y_melb = melb["Price"]

X_tr, X_te, y_tr, y_te = train_test_split(X_melb, y_melb,
                                           test_size=0.2, random_state=42)
melb_pipeline.fit(X_tr, y_tr)
score = melb_pipeline.score(X_te, y_te)
print(f"Melbourne Housing R^2: {score:.4f}")
\end{minted}

% ============================================================
\section{Exercices}

\begin{exercise}
\begin{enumerate}
  \item Construisez un pipeline \texttt{ColumnTransformer} pour Titanic avec : imputation m\'ediane + mise \`a l'\'echelle pour le num\'erique, et imputation mode + one-hot pour le cat\'egoriel. \'Evaluez par CV 5-fold.
  \item Cr\'eez un transformateur personnalis\'e \texttt{OutlierClipper} qui clippe aux 1\textsuperscript{er} et 99\textsuperscript{e} percentiles. Int\'egrez-le dans un pipeline.
  \item Construisez un pipeline complet pour Melbourne Housing qui inclut du feature engineering (prix par pi\`ece, \^age du b\^atiment). Sauvegardez-le avec \texttt{joblib} et rechargez-le pour pr\'edire.
  \item Comparez la performance d'un pipeline avec et sans pr\'etraitement sur Adult Census. Utilisez un arbre de d\'ecision.
  \item Cr\'eez un pipeline qui g\`ere simultan\'ement num\'erique, cat\'egoriel et texte avec \texttt{ColumnTransformer}. Utilisez TF-IDF pour le texte.
\end{enumerate}
\end{exercise}

% ============================================================
\section{R\'esum\'e du chapitre}

\begin{itemize}
  \item Le \texttt{Pipeline} scikit-learn cha\^ine transformations et mod\`ele en un seul objet, \'evitant la fuite de donn\'ees.
  \item \texttt{ColumnTransformer} applique des transformations diff\'erentes \`a des types de colonnes diff\'erents.
  \item Les transformateurs personnalis\'es (h\'eritant de \texttt{BaseEstimator} et \texttt{TransformerMixin}) s'int\`egrent naturellement dans les pipelines.
  \item Sauvegarder les pipelines avec \texttt{joblib} capture tous les param\`etres ajust\'es pour un d\'eploiement reproductible.
  \item Un pipeline garantit que le pr\'etraitement \`a l'entra\^inement et \`a la pr\'ediction sont identiques.
\end{itemize}
